\section{Discussion}\label{sec:discussion}
\subsection{A valid negative result, not a broken pipeline}
The central outcome of this thesis is negative: at the noise level actually measured on the
OAK-D SR PoE, injecting a fitted ToF noise model into simulation does not change pose
accuracy (Section~\ref{sec:res-inject}) and does not change simulated grasp feasibility
(Section~\ref{sec:res-simmanip}). This refutes the proposal hypothesis. The result is sound
rather than an artefact of an insensitive pipeline, because the same injector does degrade
perception once the axial term is scaled up: pose error grows only past roughly $\times10$
measured noise ($\sigma\approx9$\,mm) and manipulation collapses only at $\times20$
(Section~\ref{sec:res-ablation}). Realistic noise simply sits an order of magnitude below
the level at which it begins to matter.

\subsection{The scope of the null: benign versus harsh environments}
\label{sec:disc-scope}
The negative result is not a claim that ToF noise never matters; it is a claim about a
specific---and very common---operating regime. All of the noise measured in this work was
recorded indoors under controlled, diffuse lighting, where the OAK-D's multi-frequency
phase unwrapping keeps axial precision sub-millimetre (Section~\ref{sec:res-noise}). The
ablation (Section~\ref{sec:res-ablation}) is the other half of the story: the estimator
\emph{does} break once the axial term is scaled to roughly $\times10$ ($\sigma\approx9$\,mm),
and grasping collapses by $\times20$ ($\sigma\approx18$\,mm). Crucially, these are not
arbitrary stress levels---they are the levels that ToF cameras genuinely reach \emph{outside}
the benign indoor regime. Under strong ambient infrared such as direct sunlight, in hot or
industrial settings, in airborne particulates (dust, fog, smoke), on specular or dark
low-reflectivity surfaces, and at extended range, ToF axial error is well documented to grow
by an order of magnitude or more, driven by a reduced signal-to-background ratio, increased
shot and thermal noise, and multipath \cite{foix2011lockin,sarbolandi2015kinect}. The
practical reading of the negative result is therefore \emph{conditional on the environment}.
For an indoor or office tabletop system---the setting this thesis targets---modelling
realistic depth noise is not worth the effort, because the operating point sits an order of
magnitude below the level at which noise reaches the median-centroid estimate. For a system
meant to run outdoors, in a hot or hazardous industrial environment, or at long range, the
same analysis predicts the opposite: the operating point moves into the $\times10$--$\times20$
band where injected noise demonstrably degrades both pose accuracy
(Table~\ref{tab:pose}) and grasp success (Table~\ref{tab:grasp}), and there a faithful
ToF-noise model---together with a perception front-end designed to tolerate it---would pay
off. The ablation should thus be read not as a curiosity an order of magnitude from reality,
but as a map of the boundary between these two regimes, and it is the regime, not the noise
model alone, that decides whether depth-noise fidelity is worth modelling.

\subsection{Why the estimator absorbs realistic noise}
The pose-relevant noise on the real sensor does not live in the body of each object but in
a thin ring of edge mixed-pixels, whose flying-pixel rate reaches 53 to 99\% across the
four objects while the bulk axial precision stays below $1$\,mm
(Section~\ref{sec:res-noise}). A per-axis median centroid (Section~\ref{sec:methods-perc})
has a 50\% breakdown point \cite{huber1981robust}, so it rejects this edge population as
long as the outliers stay a minority of the blob. The injected axial Gaussian, by contrast,
perturbs the quiet body pixels that the median already estimates well, so it has almost no
leverage on the centroid. This also reframes the literature models: rather than adopting
Nguyen's quadratic range term blindly, the Phase~2 data test it and find it inapplicable to
this multi-frequency sensor \cite{nguyen2012modeling}.

\subsection{Where the gap actually lives}
If depth-noise fidelity does not govern the gap, what does? The real-versus-sim planar pose
comparison (Section~\ref{sec:res-real}) puts the end-to-end gap at $-1.4$ to $+3.6$\,mm
across all four objects, with clean and noisy simulation still indistinguishable on real
data. The residual is consistent with calibration and lighting rather than sensor
stochasticity: the extrinsic is a single direct \texttt{solvePnP} solution, factory
intrinsics are used without a dedicated calibration run, and segmentation quality varies
with illumination at the object boundary. This echoes the broader lesson of domain
randomisation work, that the factors which dominate transfer are often not the ones a given
pipeline randomises \cite{tobin2017domain}, and it locates the reality gap
\cite{jakobi1995noise} for this system in geometry and photometry rather than in the
depth-noise statistics.

\subsection{Manipulation: the real bottleneck is kinematic, and grip mechanics decide}
The manipulation results separate two error sources. In simulation, grasp feasibility is
governed purely by perceived-centre error against jaw tolerance, so the three objects with
adequate tolerance are flat from clean through $\times10$ and only the size-limited coke can
is floored. On the real arm, the dominant error is not perception but the kinematic model:
the xArm~1S over-reaches its URDF by roughly 28\%, and a global similarity-scale workaround
leaves about $14$\,mm of residual position error, more than the blue box's $\pm7.5$\,mm
capture tolerance, which confined open-loop grasping to a narrow patch. The closed-loop
visual nudge (Section~\ref{sec:methods-servo}) removes this by construction: because the
error is a difference between two points seen through the same overhead camera in the same
frame, the wrong forward kinematics and the scale factor cancel as common-mode bias, and
grasp accuracy becomes camera-bound rather than model-bound. With aiming corrected, the
per-object grasp rates (Table~\ref{tab:phase4c}) are ordered by neither the simulated
feasibility ceiling nor the measured pose error, but by the interaction of object shape with
the parallel-jaw closing force: the coke can rises from a predicted 0\% to 67\% (a firm close
on a slightly compliant, near-jaw-width can), while the golf ball falls from a favourable
65\% prediction to 20\% (a smooth sphere that rolls out of flat jaws). This is consistent
with the grasping literature's distinction between geometric reachability and grasp
stability \cite{bohg2014grasp}.

\subsection{Simulation realism and the analytic metric}
A deliberate scoping choice shapes how the simulated manipulation numbers should be read.
Gazebo Classic \cite{koenig2004gazebo} cannot reliably resolve contact grasps on small or
curved objects in this setup; finger collisions pass through and object pose is unchanged
across a jaw close. The simulated success metric is therefore analytic grasp feasibility,
not scored physical contact, and it was renamed accordingly to avoid overclaiming. This
keeps every simulated error attributable to perception, but it means the simulated rates are
an upper bound on what a physical grasp would achieve. The contact-attach plugin
\cite{link_attacher} is used only for visualisation.

\subsection{Limitations}
Several limitations bound the strength of these claims. The real perception comparison is
$n=4$ per object: in four of five scenes one object drifted into the ToF edge zone and
failed to segment, because the reliable ToF area is smaller than the RGB marker field of
view. The real manipulation evaluation covers all four objects (60 trials) but is restricted
to the near workspace ($r\le0.20$\,m); beyond that radius the loop gain falls because the
marker is mounted high and behind the jaws, a known mounting issue rather than a limit of the
method. The Phase-4A confidence intervals are Wald and collapse to zero width at 0\% and
100\%, so the Wilson intervals used in Phase~4C are the more honest choice. The noise model
omits spatial correlation and multipath, and the lateral jitter term uses the literature
value of $0.8$\,px \cite{nguyen2012modeling} rather than a separately measured one. Finally,
Gazebo Classic reached end of life in January~2025, so while the simulation results are
stable they rest on an unsupported platform. None of these affect the central negative
result, which is established independently in Phases~2 and~3 and corroborated on real data in
Phase~4B.

% =====================================================================
